% results_a.tex : Results, first four subsections (evaluation design, zero labels,
% three to ten labels, method selection). Body text only, no preamble.
% Spec: paper/manuscript/MANUSCRIPT_SPEC.md sections 4.1 to 4.4 and 11.
% Numbers: copied from paper/claims/{D_data_and_design,C7,C1,C2,C3,C5}/README.md and
% rounded with style guide 4.5 (|delta| < 1 cm two decimals, other cm values one decimal,
% correlations two decimals, percentages integer, CI ends with the digits of the printed
% point estimate). Rounding used the full-precision source rows that the READMEs cite.
% Replace with numbers.tex macros when build_numbers.py exists (QA item M-09).
% Sign: negative values use U+2212. Intervals are cell-weighted unless marked block-equal.
% Bracketed new-experiment placeholders are the registered strings of spec section 11 and
% must be replaced only by the pre-fixed interpretation sentences. The two bracketed
% missing-value markers follow spec section 12, item M-04 (median and mean without W Russia).
% Placeholders contain Korean text, so a build that keeps them needs a Korean font setup.
% Citation keys follow refs_intro.bib and refs_discussion.bib (FirstAuthorYear) and are
% listed in paper/manuscript/en/refs_results_a.bib.
% This file opens the Results section; results_b.tex continues with subsections only.

\section{Results}\label{sec:results}

\subsection{Evaluation design and validation ladder}\label{sec:results-design}

All methods were compared on one protocol in which each target region and a 100~km buffer were removed from the source data, labels were drawn from half of the target's 0.5° blocks, and errors were scored on the other half, for label counts from zero to all available rows (30 to 13,606 rows per region) (Fig.~1c,d; Table~1).
Every method was compared with two Stefan baselines that used the same target labels: the source-coefficient Stefan model, whose coefficient $E_0$ was estimated from the source cells (1.59--1.62~cm~(°C~d)$^{-1/2}$ for the main regions, whose source cells were 78--94\% Alaskan, and 1.50 for the Alaskan reference target), and the Stefan model recalibrated with the target labels by shrinkage towards $E_0$ ($\kappa = 10$).
Four-region means weight the Lena Delta, Canada, W Russia and E Russia equally and exclude Alaska, which is a held-out reference target.
The transfer comparisons were internally pre-registered, intervals are cell-weighted unless marked block-equal, and error floors used to scale effect sizes were 11.3, 14.8 and 17.9~cm in Alaska, the Lena Delta and Canada (Methods).

Averaged over five regions with equal weight, the RMSE of direct ML increased from 22.9~cm under random cell splits to 33.3~cm under region holdout (difference 10.3~cm; 95\% confidence interval (CI) 7.3 to 12.8; block-equal 6.1~cm, 4.7 to 7.6), whereas a Stefan model with the coefficient fitted by least squares to the training cells stayed between 26.8 and 27.0~cm (Fig.~1e).
The 10.3~cm increase is reported as an effect size rather than a confirmatory result, and the registered contrast between cluster holdout with a 500~km buffer and random cells was 8.2~cm (95\% CI 5.6 to 10.4; Supplementary Table S2).
Within Alaska, the median distance from test cells to the nearest training cell was 0.004~km under random splits and 53.5~km under k-fold nearest-neighbour distance matching (kNNDM) cross-validation, which mimics the 81.6~km from map cells to the nearest label\cite{Linnenbrink2024}; the RMSE of direct ML rose from 11.6 to 17.9~cm, and that of the Stefan model stayed at 14.2 to 14.5~cm (means of three seeds).
Random validation of clustered samples can overstate model accuracy\cite{Ploton2020}, and model-free, unbiased estimates of map accuracy have been argued to require probability sampling\cite{Wadoux2021}.
Our labels are a clustered non-probability sample, so the errors below are prediction errors on held-out blocks.
[XH: R1 | 결과 열람 뒤 설계, 서술(판정어 없음). L19, L20 은 등록 | 2.8 사전 고정 서술 규칙]

\subsection{Transfer without target labels}\label{sec:results-label0}

Without target labels, pseudo-labels from the Stefan model reduced RMSE relative to shuffled pseudo-labels by 1.6~cm (95\% CI 1.0 to 2.0; block-equal 0.6 to 1.5; four regions; Holm-adjusted \textit{P} = 0.001), and in a post hoc count, fewer of 30 targets lost more than 2~cm against the source-coefficient Stefan model with physics pseudo-label augmentation (2 targets) than with direct ML (18 targets) (Fig.~3b; Fig.~6a,b).
Relative to two constant placebos the reductions were 1.7 and 3.6~cm, and relative to a placebo linear in the square root of thawing degree-days 0.48~cm without labels (below 0.5~cm), with equivalence within ±0.5~cm at ten labels (registered verdict: mixed).
The 2~cm threshold was set after the transfer results had been viewed, and the 30 targets comprise seven regions and 23 sub-regions or label-expanded versions of them; counted by intervals above zero under both weightings, 14 targets had higher error with direct ML and 8 with physics pseudo-labels (post hoc computation; Supplementary Table S5).

For the five learners whose primary and common-resampling intervals agreed, direct ML had higher four-region mean error than the source-coefficient Stefan model: a random forest (2.6~cm) and four neural networks configured by leave-one-source-region-out cross-validation (2.4 to 5.2~cm) (Fig.~6c).
For the main learner, CatBoost, the difference was 2.2~cm (95\% CI 1.1 to 3.4; block-equal 1.6~cm, 0.5 to 2.7; Holm-adjusted \textit{P} = 0.023), and for two further CatBoost settings, TabPFN v2 and TabICL v2 it was 0.85 to 1.9~cm.
These five verdicts depend on the split-independence assumption, and the TabPFN v2 verdict was significant before correction only (Holm-adjusted \textit{P} = 0.117).
Direct CatBoost also had 3.2~cm higher error than the Stefan model driven by year-matched thawing degree-days (95\% CI 1.9 to 4.3; block-equal 2.0~cm, 0.9 to 3.1; auxiliary hypothesis with an unblinded anchor variant).
Residual models with a low weight ($\lambda = 0.25$) on the source-coefficient anchor were equivalent to that model within ±0.5~cm for CatBoost, TabPFN v2 and TabICL v2, a verdict that depends on the split-independence assumption.

Three results depart from this pattern.
With Alaska held out, physics pseudo-labels had 1.8~cm higher error than the source-coefficient model (95\% CI 0.4 to 2.6).
On the Tibetan Plateau, where the source coefficient is far from the regional value (Fig.~1b), direct ML had 60.0~cm lower error (95\% CI 57.1 to 62.9; an expected result from an unblinded, cross-environment comparison).
With ten labels, direct TabICL v2 had 2.9~cm lower four-region mean error (95\% CI 1.1 to 3.6; depends on the split-independence assumption) but 6.0~cm higher error in the three-region mean that includes Alaska.
[XG: R2 | 결과 열람 뒤 설계, 등록 이탈(WRAPUP 10) | 2.7 사전 고정 해석 문장 XG-1, XG-2, XG-3, 공통 문장]
We next asked what the first ten labels changed.

\subsection{Coefficient recalibration with three to ten labels}\label{sec:results-recal}

With ten target labels, recalibrating the Stefan coefficient reduced the four-region mean RMSE relative to the source coefficient by 2.5~cm (95\% CI 1.9 to 3.0; block-equal 1.9 to 3.3; Holm-adjusted \textit{P} = 0.001), with lower error in one of four regions (W Russia, 12.0~cm), higher error in Canada (2.3~cm; 95\% CI 0.7 to 2.9) and no difference established in the Lena Delta and E Russia (Fig.~2; Fig.~3a).
[MISSING: 중앙값과 러시아 서부 제외 값, build\_numbers.py 에서 lgw\_bundle.csv 지역 행으로 생성]
Adding residual ML to the recalibrated anchor changed RMSE by a further −0.18~cm (95\% CI −0.53 to −0.01; Holm-adjusted \textit{P} = 0.136), a difference significant before correction only, below 0.5~cm and dependent on the split-independence assumption.
In W Russia, 13.5~cm of the 14.0~cm reduction relative to the source coefficient with all labels came from recalibration (post hoc decomposition).

The error reduction of label-based correction (recalibration plus residual) was rank-correlated with the source-coefficient bias measured with ten labels (Spearman $\rho = 0.38$; 95\% CI 0.17 to 0.55; 28 targets; one-sided \textit{P} = 0.0002; designed after the label-grid results had been viewed), where the reduction is that of the method selected within the target labels (next section) relative to the source-coefficient model with all labels.
The sign of the bias was not related to the reduction ($\rho = -0.01$; 95\% CI −0.17 to 0.18) (Fig.~4a,b; Supplementary Table S8).
We could not establish a relation between coefficient error and the gain of ML beyond recalibration (post hoc rank correlations −0.08 to 0.39 depending on the definition; Supplementary Table S13).
[XA: R3 | 사후 분석(재현(비맹검)) | 2.1 해석 조각]

With ten or fewer labels, the residual structure had lower four-region mean error than ML that receives Stefan outputs as inputs (physics-input ML), by 2.7~cm without labels (95\% CI 1.7 to 4.3) and 3.7~cm with ten labels (95\% CI 2.8 to 5.2; Holm-adjusted \textit{P} = 0.001) (Fig.~3c).
[MISSING: 중앙값과 러시아 서부 제외 값, build\_numbers.py 에서 lgw\_bundle.csv 지역 행으로 생성]
The physics-input model of the ten-label contrast does not receive the recalibrated coefficient; in the contrast that supplies recalibrated Stefan values as inputs, the residual structure still had 2.5~cm lower error (95\% CI 1.6 to 3.3), with lower error in W and E Russia, 3.1~cm higher error in Canada (95\% CI 1.4 to 4.0; the E Russia and Canada verdicts depend on the split-independence assumption) and no difference established in the Lena Delta.

Adding physics pseudo-labels to residual ML was equivalent within ±0.5~cm with ten labels (−0.13~cm; 95\% CI −0.31 to 0.22; Holm-adjusted equivalence \textit{P} = 0.002) and increased error with 40 labels in the Lena Delta and Canada by 0.34~cm (95\% CI 0.25 to 0.49), a difference below 0.5~cm.
With three labels, the Stefan model fitted to the target labels by least squares had higher error than the recalibrated model in the Lena Delta (2.9~cm), Canada (5.4~cm) and E Russia (3.7~cm), and lower error in W Russia (5.0~cm) and on the Tibetan Plateau (81.8~cm), where the source coefficient was far from the regional value (descriptive comparison; Fig.~3d).
[XB: SI, R3 지시 문장 | 결과 열람 뒤 설계, 실행 전 등록(비맹검 부분 포함) | 2.2 사전 고정 해석 문장 XB-1, XB-2(라벨 10개 행)]

\subsection{Method selection with tens to hundreds of labels}\label{sec:results-selection}

Selecting the method by five-fold block cross-validation within the target labels reduced RMSE relative to a fixed residual recipe ($\lambda = 0.25$) by 1.3~cm with 160 labels (95\% CI 0.7 to 1.7; block-equal 0.8 to 1.8; Lena Delta and Canada; unchanged under common resampling) and was non-inferior at 40 and 160 labels (margin 0.5~cm), with both results arising in Canada; this rule was designed after the label-grid results had been viewed (Fig.~4c; Fig.~7d).
With 40 labels, the reduction of 0.87~cm (95\% CI 0.27 to 1.20) depends on the split-independence assumption.
In Canada the rule reduced RMSE by 2.1 and 2.7~cm with 40 and 160 labels; in the Lena Delta no difference was established.
With ten labels in the four main regions, non-inferiority was not established (−0.25~cm; 95\% CI −1.11 to 0.73), and the rule had 0.70~cm higher error in the Lena Delta (95\% CI 0.56 to 1.03); with all labels, non-inferiority depends on the split-independence assumption.
The rule did not increase the number of targets with lower error than the recalibrated Stefan model (3 against 4 for the fixed recipe with 160 labels, 2 against 8 with all labels; Supplementary Table S8).

In Canada, the rule chose residual ML with $\lambda = 1.0$ in 22 of 25 draws with 40 labels and in 24 of 25 with 160 labels, and no difference from the source-coefficient Stefan model was established.
With Alaska held out, the rule had higher error than the source-coefficient model with 10, 40 and 160 labels (0.91, 0.98 and 2.1~cm; descriptive curve contrasts).

At 40 and 160 labels, residual ML had 1.9 to 2.4~cm higher error than physics-input ML in the Lena Delta and Canada (the 2.4~cm contrast depends on the split-independence assumption) (Fig.~3c).
This difference arose in Canada (4.5 to 5.1~cm); no difference was established in the Lena Delta, and with Alaska held out residual ML had lower error.
At these label counts no difference was established between residual ML and a stronger recalibrated baseline (the soil-property Stefan model recalibrated with $\kappa = 10$).

Hyperparameter selection for neural networks by leave-one-source-region-out cross-validation gave learner-specific results (Supplementary Table S4): where both intervals agreed, it increased the error of direct ML for the MLP and the multi-head MLP (1.3 to 3.0~cm), no learner met the registered condition for a learner-general statement, and this selection was not compared with selection within the target labels.
[XC: R4 | 결과 열람 뒤 설계, 재사용 지역의 재검정(비맹검 부분 포함) | 2.3 사전 고정 해석 문장 XC-1, XC-2(+XC-2s), 뒤에 S-XC4 (b) 지역 손해 문장]
[XB: SI, R4 지시 문장 | 결과 열람 뒤 설계, 실행 전 등록(비맹검 부분 포함) | 2.2 사전 고정 해석 문장 XB-1, XB-2(라벨 40, 160 행, 레나·캐나다 2지역)]
